% --- Limitations (appendix placement, R92 M9) --------------------------------
% LIMITATIONS, full version. Emitted AFTER the bibliography beside the Ethics
% Statement so that the 9-page main-text content budget (Sections 1-6) is
% untouched; the main text keeps no bridge. This is the same treatment the
% ICLR 2026 author guide gives the Ethics and Reproducibility statements and
% follows the aggregate's item-12 footnote that Limitations, if not
% exempt, must fit on the main pages. Provenance: patch_plan.md A-L0, C1,
% 2026-08-21.

\section{Limitations}
\label{sec:limitations}

\paragraph{Character-unit reconstruction.}
The OpenBookQA character floor, added after the original audit flagged an evidence hole, is now recomputed from raw parquet and passes a 12/12 character/token self-test; the narrower remaining limitation is that character values are reconstructed from raw option text rather than stored per-item in the scored records. Item matching to the token leg therefore follows the shared loader and self-test, not a stored character-count column. The convention of excluding the continuation's leading space was explicit; including it yields the identical 0.363500 value.

\paragraph{One family at one depth; unmatched healing corpora.}
The prospective healed comparison has unmatched family, depth, corpus, and exposure. Table~\ref{tab:limitations-audit} places those quantities on aligned rows; the Qwen3 arm also cannot consume OLMo-2-token Dolmino, whose raw text is on neither disk. The design therefore supports neither a depth gradient nor a family contrast.

\paragraph{Regime, pooling, and the planned causal contrast.}
Tables mixing truncate-only non-OLMo arms with prune-then-heal OLMo-2 arms cannot identify a family effect; pooled results must not be read as per-benchmark verdicts (the Cho et al.\ overlap audit used arXiv v4, the January 26, 2026 camera-ready being unreachable from the audit network). The two planned contrasts are not merely unmeasured: their failed prerequisites are reported in Table~\ref{tab:limitations-audit}. A future read-out can ask only whether the healed arm has material item-level signal.

\begin{table}[H]
\centering
\small
\setlength{\tabcolsep}{4pt}
\begin{tabular}{>{\raggedright\arraybackslash}p{0.22\linewidth} >{\raggedright\arraybackslash}p{0.39\linewidth} >{\raggedright\arraybackslash}p{0.29\linewidth}}
\toprule
Limitation & Measured constraint & Consequence \\
\midrule
Contrast coverage & One non-OLMo family at one absolute depth; Llama-2, Llama-3, and \texttt{k10}/\texttt{k12}/\texttt{k14} remain confounded & The observed ladder is a cliff, not an estimable depth gradient. \\
Healing exposure & Qwen3: 5.72 epochs, 5.541B SlimPajama tokens; OLMo-2: 1.0 epoch, 31.7B Dolmino tokens & Training corpus and exposure are unmatched. \\
Relative depth & \texttt{keep8}: 22.2\% versus 25.0\% at 36 versus 32 layers & Relative depth is untested. \\
Planned causal contrasts & \texttt{H\_heal}: $-0.139$ pp, $p=0.0964$ for the required Qwen3 comparator; \texttt{H\_family}: 0/27 cells below the permutation null & Neither prerequisite holds, so additional training cannot identify the proposed contrasts. \\
\bottomrule
\end{tabular}
\caption{\textbf{The healed comparison is not factorial and does not identify a family or depth effect.} The table aligns the otherwise scattered exposure, depth, and prerequisite diagnostics; none supports a causal family claim.}
\label{tab:limitations-audit}
\end{table}



% --- ICLR 2026 page budget -------------------------------------------------
% The two files below hold material whose main-text home now carries a shorter
% summary, relocated 2026-08-16 to meet the 9-page main-text limit. Appendices
% sit after the bibliography and are unlimited. Nothing was deleted; the
% summary/full-text pairing is recorded in the distributed Appendix.
% ---------------------------------------------------------------------------
% RELOCATED (2026-08-16, ICLR 9-page budget): this is the full development of the
% four under-specifications. It is \input from sections/09_appendix.tex, AFTER the
% bibliography. The main-text summary that replaces it in line-of-sight is
% sections/03b_nulls_summary.tex. No content was deleted in the move.
\subsection{Four Under-Specifications of the Reference}
\label{app:nulls}

This appendix develops the four degrees of freedom summarised in Section~\ref{sec:nulls}, one subsection each, and carries the two tables they are read from. The phrase ``compare to a baseline'' hides choices large enough to reverse a verdict. These choices are part of the measured construct and must be printed with the score (Table~\ref{tab:conventions}).

\subsubsection{Tie convention}

For a longest-option heuristic, tied maxima can be split fractionally, resolved by the first or last option, credited whenever gold belongs to the tied set, or counted as wrong. On ARC-Challenge, \texttt{split}, \texttt{first}, \texttt{last}, and \texttt{wrong} place all six OLMo-2 arms significantly above the content floor; \texttt{credit} moves five of six significantly below it. Table~\ref{tab:reference-degrees} aligns the MMLU-Pro oracle span, the defensible default span, and the intact-model comparator; Table~\ref{tab:conventions} gives every per-convention floor. A convention that lets an oracle harvest every tie is useful as a reductio, not as a neutral default.

\subsubsection{Length unit}

``Longest'' may mean characters or continuation tokens. Table~\ref{tab:reference-degrees} reports the matched-item shifts, the OpenBookQA contrast, and the ARC-Challenge tie-rate change in aligned columns. Their scale shows why printing the convention without the unit is not reproducible specification.

The character-unit values were recomputed from raw parquet under $\mathrm{len}(\text{option text})$. Counting the continuation's leading space gives the same OpenBookQA value, 0.363500, because the shift is uniform across options. The character and token legs are item-matched by the shared loader and a 12/12 self-test.

\subsubsection{Tokenizer}

Within the token unit, the floor is tokenizer-valued. Table~\ref{tab:reference-degrees} consolidates the audited \texttt{split} and \texttt{credit} movements and the vocabulary-size comparison. This is a floor-level consequence of the known tokenizer dependence of length-based scoring \citep{oostermeijer2026length}; we make no vocabulary-size explanation because the ordering is not monotone.

This degree of freedom creates an asymmetry. A token-based content floor is a property of $(\text{dataset},\text{convention},\text{unit},\text{tokenizer})$. A character-based floor is tokenizer-free under its stated character convention. By contrast, the letter floor is a pure property of the item set's gold labels: it is invariant across all 15 cross-family arms and bit-identical across all 21 MMLU-Pro cells.

\subsubsection{What ``chance'' means when option count varies}

MMLU-Pro is described as up to ten-way, but its option count varies item to item. Table~\ref{tab:reference-degrees} places the naive and item-average chance lines beside the same observed floor, including both additive and multiplicative gaps. Both are defensible descriptions of a chance line, so both must be reported. The floor remains a dataset property; its gap to an under-specified chance line does not.

\begin{table}[!htbp]
\centering
\small
\setlength{\tabcolsep}{3.2pt}
\begin{tabular}{>{\raggedright\arraybackslash}p{0.14\linewidth}>{\raggedright\arraybackslash}p{0.22\linewidth}>{\raggedright\arraybackslash}p{0.38\linewidth}>{\raggedright\arraybackslash}p{0.16\linewidth}}
\toprule
Choice & Audited contrast & Measured sensitivity & Reporting takeaway \\
\midrule
Tie rule & MMLU-Pro, OLMo-2 tokenizer
& \texttt{wrong}: 0.125914; \texttt{credit}: 0.532164 (40.6 pp, 4.2264$\times$).  The defensible \texttt{split}/\texttt{first}/\texttt{last} span is approximately 0.5--3.1 pp; \texttt{credit} exceeds intact \texttt{content\_norm} 0.207613 by 32.5 pp.
& Print the rule; treat \texttt{credit} only as an oracle bound. \\
Length unit & Characters versus continuation tokens
& \texttt{split} shifts by at most 2.02 pp; \texttt{credit} shifts by 14.53--35.20 pp.  OpenBookQA \texttt{credit} is 0.416 versus 0.644; ARC-Challenge tied-longest rates are 18.60\% versus 50.85\%.
& The convention alone is not reproducible specification. \\
Tokenizer & Item-matched token comparisons
& \texttt{split} shifts by 0.9003 pp; \texttt{credit} shifts by up to 10.6 pp on four-way tasks and 9.26 pp on MMLU-Pro.  The 32k/152k tokenizers yield fewer MMLU-Pro ties than the 100k/128k pair.
& Print the tokenizer; vocabulary size is not a monotone explanation. \\
Chance line & MMLU-Pro, $\texttt{n\_opt}\in\{3,\ldots,10\}$
& Naive $1/10=0.100000$ versus item-average 0.110877, against floor 0.116606: $+1.661$ pp and 1.1661$\times$ versus $+0.573$ pp and 1.0517$\times$.
& Print both defensible chance definitions. \\
\bottomrule
\end{tabular}
\caption{\textbf{Four under-specifications can move a content-side reference enough to reverse its reading.} The table moves the comparable numerical contrasts out of prose while preserving their scopes. Full per-convention floors are in Table~\ref{tab:conventions}; all values come from the same audited item records.}
\label{tab:reference-degrees}
\end{table}


\subsubsection{The calibrating null must be realisable, and we initially failed this on our own largest construct}
\label{sec:legality-aware-null}

The four degrees of freedom above concern the \emph{floor}. A fifth concerns the null used to calibrate it, and our first calibration drew uniformly over letters that are not legal on every MMLU-Pro item. That null places mass \emph{outside the support of every legal labelling of the benchmark}: it is not a conservative approximation of the right experiment but a different experiment. Table~\ref{tab:null-repair-audit} records the affected item share and the unequal legal letter marginals. Their non-exchangeability raises the maximum marginal, which is precisely the quantity the floor must be compared against.

The admissible null holds the observed \texttt{n\_opt} histogram fixed --- \texttt{n\_opt} is a property of the item set, not a random variable --- and draws each item's gold letter uniformly among \emph{its own} legal letters. Table~\ref{tab:null-repair-audit} aligns the expectation and $p$-value before and after this correction with the unchanged observed floor and downstream counts. The correction changes what the floor is compared \emph{to}, never the floor itself.

\paragraph{Which rows moved, and in which direction.} We state this explicitly because a correction that a reader cannot audit for self-service is not a correction. Where \texttt{n\_opt} is constant the two nulls are the \emph{same distribution}, so no correction is possible and none is applied. The surviving constructs are therefore exactly those whose option count was never in question, while the lost row is the largest construct and the one carrying the power analysis. Where option count varies, narrower items concentrate mass on always-legal low letters and move MMLU-Pro against us; wider exceptional items move the two ARC rows slightly in our favour. Table~\ref{tab:null-repair-audit} gives the item counts, $p$-value changes, and standard-error distances side by side. Neither ARC verdict turns, but the two directions matter because they falsify the claim that the correction can only move against the authors.

\paragraph{A rounding convention with the same sign problem.} Since $p=\Pr(\hat f\ge\text{floor})$ and the floor is an exact rational $\text{count}/n$, the threshold is an integer count. Comparing against a rounded decimal can demand an \emph{extra} count, exclude the observed outcome, and bias $p$ downward toward ``survives''. We therefore evaluate every $p$ at the exact rational floor. Table~\ref{tab:null-repair-audit} lists the affected rows, the rounded and exact values, and the clean counterexamples. No verdict changes, but the MMLU-Pro result moves away from significance rather than toward it.

\begin{table}[!htbp]
\centering
\small
\setlength{\tabcolsep}{3.0pt}
\begin{tabular}{>{\raggedright\arraybackslash}p{0.17\linewidth}>{\raggedright\arraybackslash}p{0.24\linewidth}>{\raggedright\arraybackslash}p{0.27\linewidth}>{\raggedright\arraybackslash}p{0.22\linewidth}}
\toprule
Audit item & Legality-blind or rounded form & Legal or exact form & Consequence \\
\midrule
Support & 2,051/12,032 items (17.05\%) have fewer than ten options under a uniform A--J draw & Hold the observed \texttt{n\_opt} histogram fixed and draw only among each item's legal letters & The first null assigns mass outside every legal labelling. \\
Letter marginals & $\mathbb{E}[\hat m_A]=0.110877$ versus $\mathbb{E}[\hat m_J]=0.082954$ & Compare the maximum legal marginal, not exchangeable ten-letter marginals & Exchangeability understates $\mathbb{E}[\max_L\hat m_L]$. \\
MMLU-Pro calibration & $\mathbb{E}[\hat f]=0.104457$ and $p<10^{-5}$ & $\mathbb{E}[\hat f]=0.113873$ ($+0.94$ pp), $p=0.083$; observed floor $1403/12032=0.116606$ unchanged & The balanced-null claim is withdrawn; the literal floor does not move. \\
Downstream counts & --- & 15/17, the 3/12-versus-1/12 matched-standard comparison, and off-MMLU 9/85 & All remain defined against the observed floor and are unchanged. \\
Fixed option count & MMLU: 14,042 items, four-way; BoolQ: 3,270, two-way; OpenBookQA, CommonsenseQA, and PIQA: four/five/two-way & The old and legal nulls are identical; MMLU and BoolQ retain $p<10^{-5}$ & No correction is possible or applied. \\
Variable option count & ARC-Easy: 4/2,376 five-way, $0.140\to0.136$ ($\Delta p=-0.0012$, $-3.3$ SE); ARC-Challenge: 3/1,172, $0.453\to0.448$ ($-0.0052$, $-7.3$ SE) & MMLU-Pro moves in the opposite direction because narrower items concentrate mass on always-legal letters & Both directions occur; neither ARC verdict changes. \\
Rounded threshold & MMLU-Pro $p=0.078$; PIQA $p=0.658$ & Exact rational floor gives 0.083 and 0.691; five of nine rows (four constructs) are affected & No verdict changes, but MMLU-Pro moves away from 0.05. \\
Threshold examples & MMLU 0.268908 versus $3776/14042$; BoolQ 0.621713 versus $2033/3270$ & OpenBookQA is exact at $138/500$; ARC-Easy, ARC-Challenge, and CommonsenseQA round down & A stored decimal can demand an unobserved extra count. \\
\bottomrule
\end{tabular}
\caption{\textbf{Legality and exact-threshold audit for the calibrating null.} Every numerical comparison formerly embedded in the repair prose is aligned here with its measured consequence. The generated nine-row manifest in Table~\ref{tab:app-construct-nulls} is the machine-readable source for the construct-level values.}
\label{tab:null-repair-audit}
\end{table}


\paragraph{What this leaves standing.} The protocol claim is unchanged and arguably strengthened: a null must be compatible with the construct it calibrates, and the failure mode is available to any author who writes ``uniform over $k$'' for a benchmark whose $k$ is not constant. What we withdraw is the quantitative claim for one construct. We no longer assert that MMLU-Pro's floor exceeds what a balanced null could produce; at $p=0.083$ it is unresolved at $n=12032$, close enough to the conventional line that a larger item set could plausibly settle it, and we decline to describe it as either significant or comfortably inside the noise.

\begin{table}[!htbp]
\centering
\footnotesize
\setlength{\tabcolsep}{2.4pt}
\resizebox{\textwidth}{!}{%
\begin{tabular}{lrcrrrrrrl}
\toprule
Construct & $n$ & $L^\star$ & Floor & Chance & Gap (pp) & Ratio & $\mathbb{E}[\hat f]$ & $q_{95}$ & $p$ \\
\midrule
MMLU-Pro, naive & 12032 & A & 0.116606 & 0.100000 & +1.661 & 1.1661$\times$ & 0.113873 & 0.117188 & 0.083 \\
MMLU-Pro, item-avg. & 12032 & A & 0.116606 & 0.110877 & +0.573 & 1.0517$\times$ & 0.113873 & 0.117188 & 0.083 \\
MMLU & 14042 & D & 0.268908 & 0.250000 & +1.891 & 1.0756$\times$ & 0.254350 & 0.258225 & $<10^{-5}$ \\
OpenBookQA & 500 & A & 0.276000 & 0.250000 & +2.600 & 1.1040$\times$ & 0.273190 & 0.294000 & 0.384 \\
ARC-Easy & 2376 & C & 0.266414 & 0.250161 & +1.625 & 1.0650$\times$ & 0.260523 & 0.269781 & 0.136 \\
ARC-Challenge & 1172 & B & 0.265358 & 0.250156 & +1.520 & 1.0608$\times$ & 0.264984 & 0.278157 & 0.448 \\
CommonsenseQA & 1221 & B & 0.208845 & 0.200000 & +0.884 & 1.0442$\times$ & 0.214994 & 0.226863 & 0.853 \\
PIQA & 1838 & B & 0.504897 & 0.500000 & +0.490 & 1.0098$\times$ & 0.509300 & 0.522851 & 0.691 \\
BoolQ & 3270 & B & 0.621713 & 0.500000 & +12.17 & 1.2434$\times$ & 0.506972 & 0.517125 & $<10^{-5}$ \\
\bottomrule
\end{tabular}%
}
\caption{\textbf{Only MMLU and BoolQ have floors a balanced null could not plausibly produce ($p<10^{-5}$).}
All nine constructs' floors lie at or near nominal chance, but the calibrated estimator $\hat f=\max_L \hat m_L$
is a maximum over $k$ noisy marginals on a single finite item set, so a positive Gap is not by itself
evidence that the floor differs from chance. $\mathbb{E}[\hat f]$ and $q_{95}$ are the mean and 95th
percentile of $\hat f$ under the legality-aware balanced null (each item's gold letter drawn uniformly
among \emph{its own} legal letters); $p=\Pr(\hat f\ge\text{observed floor})$ under that null. Evidence
\textsf{E-CAL} is the single machine-readable source for these rows.}
\label{tab:nulls}
\end{table}

\FloatBarrier
\begin{table}[!htbp]
\centering
\small
\setlength{\tabcolsep}{4.5pt}
\begin{tabular}{llrrrrr}
\toprule
Dataset & Unit/tokenizer & \texttt{split} & \texttt{first} & \texttt{last} & \texttt{credit} & \texttt{wrong} \\
\midrule
MMLU-Pro & token, OLMo-2 & 0.193150 & 0.195894 & 0.190824 & \textbf{0.532164} & 0.125914 \\
OpenBookQA & character & 0.363500 & 0.362000 & 0.362000 & 0.416000 & 0.320000 \\
OpenBookQA & token, OLMo-2 & 0.368000 & 0.384000 & 0.360000 & \textbf{0.644000} & 0.228000 \\
ARC-Challenge & character & 0.274104 & 0.272184 & 0.274744 & 0.347270 & 0.220990 \\
ARC-Challenge & token, OLMo-2 & 0.283902 & 0.265358 & 0.296075 & \textbf{0.543515} & 0.142491 \\
Winogrande, control & character & 0.491713 & 0.492502 & 0.490923 & 0.581689 & 0.401736 \\
Winogrande, control & token, OLMo-2 & 0.501184 & 0.494081 & 0.508287 & \textbf{0.933702} & 0.068666 \\
\bottomrule
\end{tabular}
\caption{The content floor changes with tie convention and length unit. Numbers are the content-side longest-option floors recomputed per (convention, unit, tokenizer) from the per-item option records, including the raw-parquet character reconstruction and its character/token self-test (evidence \textsf{E-F}, \textsf{E-H} and the August 14 character-unit addendum; see the reproducibility statement for the artifact paths).}
\label{tab:conventions}
\end{table}


\FloatBarrier

% RELOCATED (2026-08-16, ICLR 9-page budget). Each subsection below is the full text
% of a passage whose main-text home now carries a shorter in-text summary. Nothing was
% deleted in the move; the pairing is recorded in the distributed Appendix.

\subsection{Prior art on option-count-aware chance correction}
\label{app:priorart}

Section~\ref{sec:v2} claims the within-stratum permutation and its use as a gate, and nothing else. This appendix attributes the rest, result by result. Bennett's $S$ sets $p_e=1/k$ for $k$ response categories \citep{bennett1954communications}; \citet{brennan1981kappa} recommend exactly that free-marginal alternative wherever $\kappa$'s marginal dependence misleads; formula scoring corrects for guessing as an explicit function of the per-item option count \citep{frary1988formula}; and \citet{brenner1996weightedkappa} show that a weighted $\kappa$ drifts systematically with the number of categories --- the very pathology our stratification is designed to avoid. Pooling an agreement coefficient across heterogeneous item groups is likewise established \citep{devries2008pooledkappa}. We therefore claim \emph{none} of the following: correcting for the number of options, choosing $p_e$ as a methodological decision, or noticing that $\kappa$-family statistics depend on $k$. The works above assume a single global $k$; what we are not aware of in that literature is a null in which $k$ \emph{varies item to item} ($\texttt{n\_opt}\in\{3,\dots,10\}$, eight strata here), with the permutation confined within strata so that a letter illegal for an item can never be credited to it. That is a loophole-closing construction, not a new statistic, and its measured footprint on MMLU-Pro is the $36$ items quantified in Appendix~\ref{app:ordering}.

\subsection{Why the v1/v2 ordering is not an identity}
\label{app:ordering}

Section~\ref{sec:v2} states the ordering between the two nulls as a measured property under an explicit regularity condition. This appendix gives the argument in full, and Table~\ref{tab:two-nulls} sets the two read-outs side by side.

Unstratified, $\sum_L p_{L}m_{L}\leq\max_Lm_{L}$ places the v2 null at or below the v1 best-constant floor. Our estimator, however, permutes \emph{within} \texttt{n\_opt} strata, and stratification only yields $\widehat{\mathrm{acc}}\leq\sum_s w_s\max_L m_{s,L}$, whose right-hand side dominates $f_{\mathrm{const}}=\max_L\sum_s w_s m_{s,L}$. The ordering is recovered only when the same letter maximises every stratum, which MMLU-Pro violates. Table~\ref{tab:ordering-audit} aligns the global and stratum maxima, the combinatorial gap, the non-independent empirical count, the binding cell, and a legal counterexample. The ordering becomes a theorem under a condition we can state but not assume --- that the arm's letter distribution not vary with \texttt{n\_opt}, since $p_{s,L}=p_{L}$ gives $\widehat{\mathrm{acc}}=\sum_L p_{L}f_{L}\leq f_{\mathrm{const}}$ exactly. We therefore state it as a measured property under that explicit regularity condition, not as an identity.

\begin{table}[!htbp]
\centering
\small
\setlength{\tabcolsep}{4pt}
\begin{tabular}{p{0.12\linewidth}p{0.31\linewidth}p{0.29\linewidth}p{0.17\linewidth}}
\toprule
Read-out & Question & Null & Required dependence \\
\midrule
v1, headline & Is this interface valid for comparing arms? & Best constant, e.g.\ always-A 0.116606 & Arm-independent \\
v2, companion & Does this arm's prediction vector contain item-level information? & Prediction vector permuted within \texttt{n\_opt} strata & Arm-conditional \\
\bottomrule
\end{tabular}

\vspace{4pt}
\begin{tabular}{@{}lrrr@{}}
\toprule
Constant emitter & always-A & always-E & always-J \\
v1 delta (pp) & 0.000 & $-2.111$ & $-3.807$ \\
\bottomrule
\end{tabular}
\caption{\textbf{The two nulls answer different questions, and v1 is not collapse-invariant.} The upper panel gives their estimands and required dependence; the lower panel shows that equally uninformative constant emitters receive unequal v1 deltas on MMLU-Pro. Values and definitions come from the stratified permutation-null record for the re-judged cells (evidence \textsf{E-D}; see the reproducibility statement for the artifact path).}
\label{tab:two-nulls}
\end{table}

\begin{table}[!htbp]
\centering
\small
\setlength{\tabcolsep}{3.2pt}
\begin{tabular}{>{\raggedright\arraybackslash}p{0.18\linewidth}>{\raggedright\arraybackslash}p{0.28\linewidth}>{\raggedright\arraybackslash}p{0.25\linewidth}>{\raggedright\arraybackslash}p{0.19\linewidth}}
\toprule
Audit slice & Observed structure & Bound or count & Interpretation \\
\midrule
Global versus strata & Global best constant A: $1403/12032$; stratum argmax B/E/B/E for $\texttt{n\_opt}=5/6/7/8$ & Four of eight strata, 623/12,032 items; stratified bound $1439/12032$ & Maximum possible inversion is 36 items, or 0.299 pp. \\
Empirical cells & 27 total cells & 13 clear the floor by more than 0.299 pp; remaining 14 span four base models, including five nearby healing checkpoints & The raw 27/27 ordering overstates independent evidence. \\
Binding cell & Llama-2 \texttt{keep12}, unhealed & Observed margin 0.0078 pp; attainable violation 0.085 pp; 94.8\% weight in A-argmax strata & This cell supplies the tight comparison. \\
Constructive counterexample & Legal emitter A,A,B,E,B,E,A,A & Attains $f_{\mathrm{const}}+0.299$ pp & The ordering needs a regularity condition; it is not an identity. \\
\bottomrule
\end{tabular}
\caption{\textbf{Why the v1/v2 ordering is measured rather than algebraic.} The table separates the combinatorial bound, the empirical count, the binding cell, and a legal counterexample. Exact per-cell v2 values are in Table~\ref{tab:v2-full}.}
\label{tab:ordering-audit}
\end{table}


\FloatBarrier

\subsection{Designated damaged cells, multiplicity, and power}
\label{app:designated}

\paragraph{Designated damaged cells.} Throughout, \emph{designated damaged} means an arm that was structurally pruned and is reported as a damaged rung of its family; the designation is fixed by construction, not by measured score. The set comprises the named OLMo-2 prune-then-heal arms and the same four truncation rungs for each non-OLMo family. Tables~\ref{tab:mmlupro} and~\ref{tab:offmmlu-rollup} state the resulting denominators, while Winogrande remains a control and enters none. An earlier roll-up silently omitted \texttt{shortgpt16} and \texttt{keep14}; their floor deltas and tests are now printed in Table~\ref{tab:mmlupro}, and both arms are included everywhere. The omission was not defensible because the matching retained depth was counted as damaged in the other families. The checker \texttt{gate\_designated\_denominator.py} now asserts that the declared arm set equals every roll-up set, that depth is treated consistently across families, and that each denominator equals arms times benchmarks.

\paragraph{Multiplicity.} We report per-cell decisions without a family-wise correction, so the headline totals are descriptive counts rather than simultaneously valid claims. This is a deliberate limitation: the cells share items, nest arms, and share a null, making them neither independent nor exchangeable and leaving no defensible off-the-shelf family definition. Readers should treat each borderline result as a per-cell decision. The conclusions rest on the aggregate direction and the regime-and-depth concentration printed in Table~\ref{tab:offmmlu-rollup}, not on one cell crossing a threshold. This remains weaker than multiplicity control, but the concentration is load-bearing: an indiscriminate test would not place all rejections at one end of the damage ladder.

Table~\ref{tab:power} is the power disclosure that Section~\ref{sec:setup} requires before any null result on the five smaller benchmarks is interpreted.

\begin{table}[!htbp]
\centering
\small
\setlength{\tabcolsep}{4.5pt}
\begin{tabular}{lrrrr}
\toprule
Task & $n$ & \texttt{keep8} $\Delta$ & CI95 half-width & Detect $-1.389$ pp? \\
\midrule
MMLU & 14042 & $-1.389$ & 1.154 & reference \\
ARC-Easy & 2376 & $-0.800$ & 1.305 & yes, borderline \\
Winogrande, control & 1267 & $+0.552$ & 1.184 & yes \\
PIQA & 1838 & $+2.503$ & 2.775 & no \\
CommonsenseQA & 1221 & $-1.065$ & 3.399 & no \\
ARC-Challenge & 1172 & $-0.939$ & 3.882 & no \\
OpenBookQA & 500 & $-1.800$ & 6.400 & no \\
\bottomrule
\end{tabular}
\caption{\textbf{Mandatory achieved-precision disclosure (with one true power reference) for small-benchmark null results.} What is reported per cell is the paired-bootstrap CI95 half-width --- \emph{achieved precision}, not classical power: only the MMLU row satisfies the standard definition (the probability that the reference cell's interval would exclude zero given its observed effect, obtained from plug-in simulation), while the remaining rows say whether their achieved half-width would resolve the reference effect. Numbers are the paired-bootstrap half-widths recorded per cell in the five-benchmark letter/content null records (evidence \textsf{E-F}; see the reproducibility statement for the artifact path).}
\label{tab:power}
\end{table}


\FloatBarrier

\paragraph{Power is part of construct validity.} Small benchmarks produce apparently reassuring at-floor results even when observed effects are large, so any interpretation of a null result must include the achieved interval. MMLU-Pro changes the evidential status: OLMo-2 \texttt{keep8}'s interval excludes an MMLU-sized below-floor effect, making this a powered non-replication, not an absence of evidence. The ladder remains a cliff, not a gradient; no depth curve should be fitted to the sparse rungs.

\subsection{Multiplicity in the two-reference comparison}
\label{app:analysis-mult}

Two qualifications belong with the two-reference comparison of Section~\ref{sec:flips}. First, one instrument does not transfer: the floor comparison also carries an exact McNemar test, which has no analogue here because $\mathrm{mean}(1/\texttt{n\_opt})$ is not a binary predictor and generates no discordant pairs. Second, these remain per-cell decisions under a shared-item dependence structure. Table~\ref{tab:headline-reference} now prints the raw matched-standard comparison and the Benjamini--Hochberg/Bonferroni result on separate rows; after correction, neither reference retains a cell, so the joint comparison is undefined rather than a corrected version of the raw ratio. We therefore treat the raw counts as descriptive evidence that the chance side is not uniformly null, not as simultaneously valid per-cell claims. The independent-binomial tail used by an earlier draft contradicted this dependence structure and has been removed.

\needspace{4\baselineskip}
\subsection{Constant emitters and the off-MMLU replication}
\label{app:offmmlu}

Table~\ref{tab:offmmlu-cells} sets all 60 non-OLMo designated damaged cells against both deterministic references, with their intervals.

% GENERATED 2026-08-23 by C1 r47 (RULING_R148 P0-4) from the frozen per-cell record
% paper/evidence/second_mc_benchmark_crossfamily/gate2_crossfamily_nulls.json.
% 17 arms x 5 letter constructs; chance-side and floor-side deltas with paired-bootstrap
% half-widths and two-sided mid-p; OLMo-2 prune-then-heal rows are prose-cited from the
% same record's off-MMLU replication summary (sec. app:offmmlu).
{\footnotesize
\setlength{\tabcolsep}{4pt}
\begin{longtable}{@{}llrrrrrrl@{}}
\caption{\textbf{Chance-side and floor-side per-cell read-out for the 60 non-OLMo designated damaged cells.}
Each row gives the truncated arm's letter accuracy, deltas against the item-average chance line
$\mathrm{mean}(1/\texttt{n\_opt})$ and against the construct's best-constant floor, the paired-bootstrap
CI half-width (identical for both deterministic references), and the two-sided mid-$p$ against the floor;
the floor-side columns reproduce the values aggregated in Section~\ref{sec:flips} (0/60 above floor,
7/60 significantly below). $n$ per construct as in Table~\ref{tab:nulls}; Winogrande is the sole
negative control and enters none of these denominators. The 25 OLMo-2 prune-then-heal cells of the same grid are summarised in Appendix~\ref{app:offmmlu}: five \texttt{shortgpt16} cells above floor ($+11.21$ to $+49.66$ points), four \texttt{keep14} cells above floor ($+6.47$ to $+18.69$, PIQA at the floor), and 0/15 above floor for the remaining $\leq 12$-layer rungs. $^{\ddagger}$Literal constant emitter on the optimal constant with the floor on the same observed labels; the bootstrap supplies no spread and mid-$p$ is exactly $1.0000$. Evidence \textsf{E-B}.}
\label{tab:offmmlu-cells}\\
\toprule
Family & Rung (cell) & Acc. & $\Delta_{\rm chance}$ & hw & $p$ & $\Delta_{\rm floor}$ & hw & Floor verdict \\
\midrule
\endfirsthead
\multicolumn{9}{@{}l}{\emph{Table~\ref{tab:offmmlu-cells} (continued)}}\\
\toprule
Family & Rung (cell) & Acc. & $\Delta_{\rm chance}$ & hw & $p$ & $\Delta_{\rm floor}$ & hw & Floor verdict \\
\midrule
\endhead
\midrule
\multicolumn{9}{r@{}}{\emph{continued on next page}}\\
\endfoot
\bottomrule
\endlastfoot
Llama-2 & \texttt{k14} / ARC-C & 0.226109 & -2.405 & 3.925 & 0.0499 & -3.925 & 3.925 & below floor \\
Llama-2 & \texttt{k14} / ARC-E & 0.251684 & +0.152 & 2.841 & 0.3214 & -1.473 & 2.841 & at floor \\
Llama-2 & \texttt{k14} / OBQA & 0.272000 & +2.200 & 0.800 & 0.3208 & -0.400 & 0.800 & at floor \\
Llama-2 & \texttt{k14} / CSQA & 0.196560 & -0.344 & 3.522 & 0.5002 & -1.229 & 3.522 & at floor \\
Llama-2 & \texttt{k14} / PIQA & 0.493471 & -0.653 & 4.543 & 0.6095 & -1.143 & 4.543 & at floor \\
\midrule
Llama-2 & \texttt{k12} / ARC-C & 0.226109 & -2.405 & 3.925 & 0.0499 & -3.925 & 3.925 & below floor \\
Llama-2 & \texttt{k12} / ARC-E & 0.250842 & +0.068 & 2.841 & 0.2928 & -1.557 & 2.841 & at floor \\
Llama-2 & \texttt{k12} / OBQA & 0.276000 & +2.600 & 0.000 & 1.0000 & +0.000 & 0.000 & at floor $^{\ddagger}$ \\
Llama-2 & \texttt{k12} / CSQA & 0.195741 & -0.426 & 3.522 & 0.4724 & -1.310 & 3.522 & at floor \\
Llama-2 & \texttt{k12} / PIQA & 0.494559 & -0.544 & 4.598 & 0.6504 & -1.034 & 4.598 & at floor \\
\midrule
Llama-2 & \texttt{k10} / ARC-C & 0.221843 & -2.831 & 3.754 & 0.0230 & -4.352 & 3.754 & below floor \\
Llama-2 & \texttt{k10} / ARC-E & 0.243266 & -0.690 & 2.694 & 0.0960 & -2.315 & 2.694 & at floor \\
Llama-2 & \texttt{k10} / OBQA & 0.260000 & +1.000 & 3.902 & 0.4308 & -1.600 & 3.902 & at floor \\
Llama-2 & \texttt{k10} / CSQA & 0.197379 & -0.262 & 3.481 & 0.5192 & -1.147 & 3.481 & at floor \\
Llama-2 & \texttt{k10} / PIQA & 0.480413 & -1.959 & 3.863 & 0.2083 & -2.448 & 3.863 & at floor \\
\midrule
Llama-2 & \texttt{k8} / ARC-C & 0.230375 & -1.978 & 3.541 & 0.0520 & -3.498 & 3.541 & at floor \\
Llama-2 & \texttt{k8} / ARC-E & 0.263047 & +1.289 & 2.441 & 0.7845 & -0.337 & 2.441 & at floor \\
Llama-2 & \texttt{k8} / OBQA & 0.250000 & +0.000 & 5.000 & 0.3217 & -2.600 & 5.000 & at floor \\
Llama-2 & \texttt{k8} / CSQA & 0.203931 & +0.393 & 2.088 & 0.6317 & -0.491 & 2.088 & at floor \\
Llama-2 & \texttt{k8} / PIQA & 0.498368 & -0.163 & 3.455 & 0.7039 & -0.653 & 3.455 & at floor \\
\midrule
Llama-3 & \texttt{k14} / ARC-C & 0.262799 & +1.264 & 2.517 & 0.8553 & -0.256 & 2.517 & at floor \\
Llama-3 & \texttt{k14} / ARC-E & 0.267256 & +1.709 & 2.588 & 0.9620 & +0.084 & 2.588 & at floor \\
Llama-3 & \texttt{k14} / OBQA & 0.264000 & +1.400 & 6.400 & 0.7279 & -1.200 & 6.400 & at floor \\
Llama-3 & \texttt{k14} / CSQA & 0.213759 & +1.376 & 1.065 & 0.3803 & +0.491 & 1.065 & at floor \\
Llama-3 & \texttt{k14} / PIQA & 0.494015 & -0.598 & 1.360 & 0.1153 & -1.088 & 1.360 & at floor \\
\midrule
Llama-3 & \texttt{k12} / ARC-C & 0.258532 & +0.838 & 2.730 & 0.6337 & -0.683 & 2.730 & at floor \\
Llama-3 & \texttt{k12} / ARC-E & 0.250000 & -0.016 & 2.609 & 0.2095 & -1.641 & 2.609 & at floor \\
Llama-3 & \texttt{k12} / OBQA & 0.232000 & -1.800 & 6.200 & 0.1729 & -4.400 & 6.200 & at floor \\
Llama-3 & \texttt{k12} / CSQA & 0.218673 & +1.867 & 1.556 & 0.2345 & +0.983 & 1.556 & at floor \\
Llama-3 & \texttt{k12} / PIQA & 0.491295 & -0.871 & 0.843 & 0.0015 & -1.360 & 0.843 & below floor \\
\midrule
Llama-3 & \texttt{k10} / ARC-C & 0.241468 & -0.869 & 4.053 & 0.2512 & -2.389 & 4.053 & at floor \\
Llama-3 & \texttt{k10} / ARC-E & 0.236953 & -1.321 & 2.862 & 0.0438 & -2.946 & 2.862 & below floor \\
Llama-3 & \texttt{k10} / OBQA & 0.202000 & -4.800 & 6.000 & 0.0145 & -7.400 & 6.000 & below floor \\
Llama-3 & \texttt{k10} / CSQA & 0.208845 & +0.885 & 3.563 & 0.9999 & +0.000 & 3.563 & at floor \\
Llama-3 & \texttt{k10} / PIQA & 0.479869 & -2.013 & 3.863 & 0.2071 & -2.503 & 3.863 & at floor \\
\midrule
Llama-3 & \texttt{k8} / ARC-C & 0.217577 & -3.258 & 3.925 & 0.0144 & -4.778 & 3.925 & below floor \\
Llama-3 & \texttt{k8} / ARC-E & 0.239057 & -1.110 & 2.841 & 0.0633 & -2.736 & 2.841 & at floor \\
Llama-3 & \texttt{k8} / OBQA & 0.272000 & +2.200 & 1.400 & 0.5997 & -0.400 & 1.400 & at floor \\
Llama-3 & \texttt{k8} / CSQA & 0.203931 & +0.393 & 3.563 & 0.7835 & -0.491 & 3.563 & at floor \\
Llama-3 & \texttt{k8} / PIQA & 0.490751 & -0.925 & 4.271 & 0.5083 & -1.415 & 4.271 & at floor \\
\midrule
Qwen3 & \texttt{k14} / ARC-C & 0.226962 & -2.319 & 3.925 & 0.0550 & -3.840 & 3.925 & at floor \\
Qwen3 & \texttt{k14} / ARC-E & 0.250421 & +0.026 & 2.820 & 0.2801 & -1.599 & 2.820 & at floor \\
Qwen3 & \texttt{k14} / OBQA & 0.276000 & +2.600 & 0.600 & 0.9831 & +0.000 & 0.600 & at floor \\
Qwen3 & \texttt{k14} / CSQA & 0.195741 & -0.426 & 3.563 & 0.4675 & -1.310 & 3.563 & at floor \\
Qwen3 & \texttt{k14} / PIQA & 0.498368 & -0.163 & 4.597 & 0.7677 & -0.653 & 4.597 & at floor \\
\midrule
Qwen3 & \texttt{k12} / ARC-C & 0.226962 & -2.319 & 3.925 & 0.0550 & -3.840 & 3.925 & at floor \\
Qwen3 & \texttt{k12} / ARC-E & 0.250421 & +0.026 & 2.841 & 0.2811 & -1.599 & 2.841 & at floor \\
Qwen3 & \texttt{k12} / OBQA & 0.276000 & +2.600 & 0.600 & 0.9992 & +0.000 & 0.600 & at floor \\
Qwen3 & \texttt{k12} / CSQA & 0.192465 & -0.753 & 3.563 & 0.3700 & -1.638 & 3.563 & at floor \\
Qwen3 & \texttt{k12} / PIQA & 0.495647 & -0.435 & 4.625 & 0.6862 & -0.925 & 4.625 & at floor \\
\midrule
Qwen3 & \texttt{k10} / ARC-C & 0.226962 & -2.319 & 3.925 & 0.0550 & -3.840 & 3.925 & at floor \\
Qwen3 & \texttt{k10} / ARC-E & 0.250421 & +0.026 & 2.841 & 0.2795 & -1.599 & 2.841 & at floor \\
Qwen3 & \texttt{k10} / OBQA & 0.276000 & +2.600 & 0.000 & 1.0000 & +0.000 & 0.000 & at floor $^{\ddagger}$ \\
Qwen3 & \texttt{k10} / CSQA & 0.198198 & -0.180 & 3.645 & 0.5496 & -1.065 & 3.645 & at floor \\
Qwen3 & \texttt{k10} / PIQA & 0.494559 & -0.544 & 4.325 & 0.6356 & -1.034 & 4.325 & at floor \\
\midrule
Qwen3 & \texttt{k8} / ARC-C & 0.247440 & -0.272 & 2.432 & 0.1511 & -1.792 & 2.432 & at floor \\
Qwen3 & \texttt{k8} / ARC-E & 0.255051 & +0.489 & 2.883 & 0.4315 & -1.136 & 2.883 & at floor \\
Qwen3 & \texttt{k8} / OBQA & 0.260000 & +1.000 & 6.002 & 0.6105 & -1.600 & 6.002 & at floor \\
Qwen3 & \texttt{k8} / CSQA & 0.192465 & -0.753 & 3.522 & 0.3537 & -1.638 & 3.522 & at floor \\
Qwen3 & \texttt{k8} / PIQA & 0.505441 & +0.544 & 1.632 & 0.9558 & +0.054 & 1.632 & at floor \\
\end{longtable}
}


\begin{table}[!htbp]
\centering
\small
\setlength{\tabcolsep}{3.2pt}
\begin{tabular}{>{\raggedright\arraybackslash}p{0.19\linewidth}>{\raggedright\arraybackslash}p{0.20\linewidth}>{\raggedright\arraybackslash}p{0.31\linewidth}>{\raggedright\arraybackslash}p{0.20\linewidth}}
\toprule
Slice & Denominator/read-out & Measured result & Takeaway \\
\midrule
All designated cells & 85 cells & 45/85 above chance; 13/85 point-estimate above floor; 9/85 interval-level above floor & Every interval clear is an OLMo-2 healed arm. \\
Literal constants & Nine damaged cells plus seven control cells & Sixteen constant emitters total.  Two OpenBookQA cells score 0.276000, $\Delta=0$, CI95 $[0,0]$, yet sit 2.6 pp above chance. & Chance can credit an empty emitter; the floor does not. \\
OLMo-2, at least 14 layers & Ten cells across \texttt{shortgpt16}/\allowbreak\texttt{keep14} & Five \texttt{shortgpt16} clears ($+11.21$ to $+49.66$ pp) and four \texttt{keep14} clears ($+6.47$ to $+18.69$ pp); keep14 PIQA is at floor & Nine of ten interval-clear; a healed-ladder property. \\
Other regimes & 15 shallower OLMo-2 cells; 60 truncate-only non-OLMo cells & 0/15 and 0/60 interval-level clears.  Non-OLMo has 4/60 point-estimate above and 5/60 exactly at floor. & Equal depth does not reproduce the healed-ladder behavior. \\
Non-OLMo precision & 60 cells & 25/60 above chance; 7/60 significantly below floor; 52/60 underpowered for the MMLU effect.  ARC-Challenge: effect $-3.840$ pp, half-width 3.92; MMLU: $-3.603$, 1.18. & Most small-benchmark nulls remain non-observations. \\
Boundary separation & Nine clears versus 76 remaining cells & Smallest interval-clear margin $+6.47$ pp; largest remaining point estimate $+2.61$ pp & The two groups do not overlap on the reported margins. \\
\bottomrule
\end{tabular}
\caption{\textbf{Off-MMLU roll-up by regime, depth, and achieved precision.} This table consolidates the comparable counts and margins from the 60-cell Table~\ref{tab:offmmlu-cells} and the paired-bootstrap evidence record; Winogrande remains a control and enters no damaged-cell denominator.}
\label{tab:offmmlu-rollup}
\end{table}


The most vivid small-benchmark cases are literal constant emitters. Table~\ref{tab:offmmlu-rollup} separates damaged cells from the negative control and aligns the OpenBookQA constant-emitter score, floor delta, interval, and chance gap. The Winogrande letter-side control floor is $0.504341$ (evidence \textsf{second\_mc\_benchmark/gate2\_letter\_content\_nulls.json}; arm-invariant across the studied arms), and the control still enters no damaged-cell denominator. The empty emitters illustrate the central discrepancy: chance credits them, while the optimal-constant reference does not.

The full off-MMLU roll-up is reported in Table~\ref{tab:offmmlu-rollup}: chance, point-floor, and interval-floor counts are shown on one row; healed, shallower, and truncate-only regimes are separated; the margin gap and the achieved-precision contrast are explicit. The interval-level clears are concentrated at the top of the OLMo-2 healed ladder, while equal or shallower retained depths in the other regimes do not reproduce that behaviour. The non-overlapping reported margins support that concentration, whereas the achieved half-widths explain why most small-benchmark nulls remain non-observations rather than powered replications. Regime and family remain confounded here as everywhere in this paper: OLMo-2 is the only prune-then-heal family, so the depth threshold is a property of that ladder and must not be read as a family contrast.

\subsection{Auditing the audit: the three defects in full}
\label{app:audit}

Significance alone would re-import the defect because a tiny effect can be detected at MMLU-Pro scale while one letter dominates the likelihood argmax. Table~\ref{tab:v2-alternatives} aligns the rejected own-modal, content-interface, entropy, and modal-share diagnostics with their same-cell comparators: entropy and modal share are descriptive rather than competence measures.

\paragraph{Full precision falsifies the numerical-tie mechanism.}
The full-precision row of Table~\ref{tab:diagnostic-contrasts} removes every exact tie and reshuffles letter decisions without recovering accuracy, and the arm remains below its floor: precision does not create item-level information. The analysis likewise repaired three implementation defects (doubled-tail bootstrap, tokenizer-specific left truncation, OOM partial result), none of which changed a substantive verdict (Table~\ref{tab:integrity}).

The analysis uncovered and repaired three implementation defects without changing a substantive verdict. Table~\ref{tab:integrity} aligns each defect with its consequence, repair, and measured impact: the bootstrap repair makes the probability bound structural; the tokenizer-specific truncation is now a hard per-shard assertion; and the cache-disabled teacher-forced pass recovers the previously unavailable intact cell. All reported cross-family MMLU-Pro results use the corrected launch.


\subsection{Per-cell MMLU-Pro read-out and the v2 re-sort}
\label{app:mmlupro-and-resort}

Table~\ref{tab:mmlupro} reports all 17 designated damaged MMLU-Pro cells.
Table~\ref{tab:v2-resort} isolates the five verdict changes, and
Table~\ref{tab:v2-full} gives the complete 27-cell v2 family.

\begin{table}[!htbp]
\centering
\small
\setlength{\tabcolsep}{4pt}
\begin{tabular}{@{}llll@{}}
\toprule
Cell & $\Delta_{\rm lit}$ & $\Delta_{\rm cal}$ & Interval verdict \\
\midrule
OLMo-2 keep8 & below & above & unchanged (at floor) \\
Llama-2 k12 & below & above & unchanged (at floor) \\
Llama-3 k14 & below & above & unchanged (at floor) \\
Qwen3 k12 & below & above & unchanged (at floor) \\
\midrule
OLMo-2 keep14 & $+0.324$ & $+0.597$ & at floor (straddles lit.\ floor) \\
Qwen3 k14 & clears & clears & clears (both refs, both stds) \\
\bottomrule
\end{tabular}
\caption{\textbf{Reference anatomy for the five cells separating the literal floor from the frozen calibrated expectation, plus one clear-both comparator.} Sign-flip cells reverse the sign of $\Delta$ between the two references without moving their interval verdict; \texttt{OLMo-2 keep14} and \texttt{Qwen3 k14} differ in kind as the table indicates. Exact per-cell margins in Table~\ref{tab:mmlupro}.}
\label{tab:reference-anatomy}
\end{table}


% v2-per-cell read-out table (C1 patch B-L1, 2026-08-21).
% Adds: (i) a point-estimate column against $\mathbb{E}[\hat f]=0.113873$, the
% legality-aware balanced-null expectation of the winning-constant score;
% (ii) an explicit propagation caveat: the paired bootstrap is taken at the
% observed winning letter, so the reference is treated as known. The
% calibration record measures the reference's own sampling distribution at
% sd=$0.2446$pp over $200{,}000$ multinomial draws at the observed gold
% marginal; the per-cell decisions below are therefore conditional on the
% observed floor. Where the reviewer-derived "5 clearing, 12 at-or-below"
% appears, the shipped rows support 2/17 significantly
% above either reference by the paired-bootstrap criterion alone, and the
% reviewer's point-estimate tally 9/17 above $\mathbb{E}[\hat f]$ versus 5/17
% above the literal floor is what the paired bootstrap does not settle --
% it is the reason this table now prints both references.

\begin{table}[H]
\centering
\footnotesize
\setlength{\tabcolsep}{2.4pt}
\resizebox{\textwidth}{!}{%
\begin{tabular}{@{}llrrrrrllrrr@{}}
\toprule
Family & Arm/regime & Acc. & $\Delta_{\rm lit}$ & $\Delta_{\rm cal}$ & hw$_{\rm floor}$ & $p_{\rm floor}$ & Lit. verdict & Cal.\ label & $\Delta_{\rm chance}^{\dagger}$ & hw$_{\rm chance}$ & $p_{\rm chance}$ \\
\midrule
OLMo-2 & \texttt{shortgpt16}, heal & 0.153341 & $+3.674$ & $+3.947$ & 0.702 & 0.0001 & above floor & above floor (pt-est) & $+4.246$ & --- & --- \\
OLMo-2 & \texttt{keep14}, heal & 0.119847 & $+0.324$ & $+0.597$ & 0.636 & 0.3234 & at floor & above floor (pt-est) & $+0.897$ & --- & --- \\
OLMo-2 & \texttt{keep12}, heal & 0.113115 & $-0.349$ & $-0.076$ & 0.765 & 0.3805 & at floor & untested (${-0.076}$) & $+0.224$ & --- & --- \\
OLMo-2 & \texttt{keep10}, heal & 0.112450 & $-0.416$ & $-0.142$ & 0.727 & 0.2662 & at floor & at/below floor & $+0.157$ & --- & --- \\
OLMo-2 & \texttt{keep8}, heal & 0.115442 & $-0.116$ & $+0.157$ & 0.582 & 0.7118 & at floor & above floor (pt-est) & $+0.456$ & --- & --- \\
Llama-2 & \texttt{k14}, truncation & 0.116772 & $+0.017$ & $+0.290$ & 0.083 & 0.7072 & at floor & above floor (pt-est) & $+0.589$ & 0.577 & 0.0456 \\
Llama-2 & \texttt{k12}, truncation & 0.115858 & $-0.075$ & $+0.198$ & 0.083 & 0.0693 & at floor & above floor (pt-est) & $+0.498$ & 0.575 & 0.0882 \\
Llama-2 & \texttt{k10}, truncation & 0.113198 & $-0.341$ & $-0.067$ & 0.416 & 0.1132 & at floor & at/below floor & $+0.232$ & 0.572 & 0.4210 \\
Llama-2 & \texttt{k8}, truncation & 0.107463 & $-0.914$ & $-0.641$ & 0.736 & 0.0168 & below floor & at/below floor & $-0.341$ & 0.548 & 0.2308 \\
Llama-3 & \texttt{k14}, truncation & 0.114694 & $-0.191$ & $+0.082$ & 0.852 & 0.6631 & at floor & above floor (pt-est) & $+0.382$ & 0.569 & 0.1802 \\
Llama-3 & \texttt{k12}, truncation & 0.111868 & $-0.474$ & $-0.201$ & 0.852 & 0.2847 & at floor & at/below floor & $+0.099$ & 0.566 & 0.7250 \\
Llama-3 & \texttt{k10}, truncation & 0.111951 & $-0.465$ & $-0.192$ & 0.827 & 0.2885 & at floor & at/below floor & $+0.107$ & 0.568 & 0.6938 \\
Llama-3 & \texttt{k8}, truncation & 0.113531 & $-0.308$ & $-0.034$ & 0.407 & 0.1296 & at floor & at/below floor & $+0.265$ & 0.563 & 0.3714 \\
Qwen3 & \texttt{k14}, truncation & 0.118933 & $+0.233$ & $+0.506$ & 0.191 & 0.0192 & above floor & above floor (pt-est) & $+0.806$ & 0.577 & 0.0062 \\
Qwen3 & \texttt{k12}, truncation & 0.115027 & $-0.158$ & $+0.115$ & 0.428 & 0.4645 & at floor & above floor (pt-est) & $+0.415$ & 0.579 & 0.1546 \\
Qwen3 & \texttt{k10}, truncation & 0.118185 & $+0.158$ & $+0.431$ & 0.233 & 0.1872 & at floor & above floor (pt-est) & $+0.731$ & 0.575 & 0.0126 \\
Qwen3 & \texttt{k8}, truncation & 0.107796 & $-0.881$ & $-0.608$ & 0.819 & 0.0362 & below floor & at/below floor & $-0.308$ & 0.547 & 0.2672 \\
\bottomrule
\end{tabular}}%
\caption{All 17 designated damaged MMLU-Pro cells against the observed literal
floor and a frozen calibrated reference. Deltas and confidence-interval
half-widths are percentage points. The paired-bootstrap $p$ freezes the observed
winning letter; the calibrated column is therefore a sensitivity surface, not a
replacement verdict, and is labelled accordingly: the nine cells whose $\Delta_{\rm cal}$
point estimate is positive print `above floor (pt-est)', the cell whose $\Delta_{\rm cal}$ is
negative but not interval-resolved prints the explicit `untested' marker, and the
remaining cells print `at/below floor'. These are reporting labels from the point
estimate, not separate tests against the calibrated reference.
$^{\dagger}\Delta_{\rm chance}$ uses the item-average convention
$\mathrm{mean}(1/\texttt{n\_opt})=0.110877$; its chance-side
\mbox{hw$_{\rm chance}$} and \mbox{$p_{\rm chance}$} columns are transcribed from the
same 10{,}000-resample run (evidence \textsf{s2\_03}, reference recomputed inside every
resample) and differ from the floor-side columns cell-by-cell --- the two sides are
\emph{separate} intervals, not shared. The five OLMo-2 cells' chance bootstrap was not
part of that run (columns print ---). OLMo-2 is prune-then-heal; the other families are
truncate-only. See the reproducibility statement.}
\label{tab:mmlupro}
\end{table}


Together the tables support ``Wrong-null verdict flips at MMLU-scale power'' and
``V2 re-sorts the 27-cell read-out.''  The Benjamini--Hochberg column in
Table~\ref{tab:v2-full} separates the surviving downgrade from the
non-surviving upgrade; all values are unchanged from the Results and Analysis
section.

\FloatBarrier

\begin{table}[!htbp]
\centering
\scriptsize
\setlength{\tabcolsep}{2.8pt}
\begin{tabular}{lrrlrrl}
\toprule
Cell & v1 $\Delta$ & v1 $p$ & v1 verdict & $\Delta_{\mathrm{perm}}$ & v2 $p$ & v2 verdict \\
\midrule
Qwen3 \texttt{k8}, unhealed & $-0.881$ & 0.0362 & below floor & $-0.139$ & 0.0964 & no item-level signal \\
Llama-2 \texttt{k8}, unhealed & $-0.914$ & 0.0168 & below floor & $-0.416$ & 0.1002 & no item-level signal \\
Qwen3 \texttt{k14}, unhealed & $+0.233$ & 0.0192 & above floor & $+0.267$ & 0.0066 & trace signal \\
OLMo-2 \texttt{keep14}, healed & $+0.324$ & 0.3234 & at floor & $+0.608$ & 0.0172 & trace signal \\
Llama-2 intact & $+1.538$ & 0.0001 & above floor & $+1.959$ & 0.0001 & anchor blocked, recovery 0.0545 \\
\bottomrule
\end{tabular}
\caption{V2 re-sorts rather than uniformly deflates the read-out. All deltas are percentage points. Numbers are the paired v1-floor and v2-permutation records for the same cells (evidence \textsf{E-D}; see the reproducibility statement for the artifact path).}
\label{tab:v2-resort}
\end{table}


\FloatBarrier

\subsection{Complete read-out v2 table}

Table~\ref{tab:v2-full} reports all 27 cells used in the post-hoc v2 re-analysis. the capacity-below-resolution and degenerate-zero-capacity integrity flags fire on no cell; all 27 pass the integrity gate; the permutation and bootstrap decisions disagree on 0/27 cells at $\alpha=0.05$. The largest absolute difference between their $p$-values is 0.0371, away from the decision boundary. Stratification is nevertheless load-bearing: for unhealed Qwen3 \texttt{k8}, the stratified statistic differs from the unstratified one by $-1.271$ points because E is not legal on every item.

\begin{table}[H]
\centering
\footnotesize
\setlength{\tabcolsep}{1.8pt}
\begin{tabular}{llrrrrrrl}
\toprule
Family & Cell & Acc. & Null & $\Delta$ & CI hw & $p$ & $q_{\mathrm{BH}}$ & Verdict \\
\midrule
Qwen3 & heal@5000 & 0.115276 & 0.115993 & $-0.072$ & 0.231 & 0.5412 & 0.7306 & none \\
Qwen3 & heal@5500 & 0.115691 & 0.115654 & $+0.004$ & 0.316 & 0.9810 & 0.9968 & none \\
Qwen3 & heal@6000 & 0.114860 & 0.115857 & $-0.100$ & 0.281 & 0.4932 & 0.7306 & none \\
Qwen3 & heal@6500 & 0.114943 & 0.115864 & $-0.092$ & 0.287 & 0.5204 & 0.7306 & none \\
Qwen3 & heal@7000 & 0.115775 & 0.115995 & $-0.022$ & 0.241 & 0.8522 & 0.9839 & none \\
OLMo-2 & \texttt{keep8}@45000 & 0.117271 & 0.114046 & $+0.322$ & 0.370 & 0.0856 & 0.2185 & none \\
OLMo-2 & \texttt{keep8}@121000 & 0.115442 & 0.113152 & $+0.229$ & 0.446 & 0.3122 & 0.5620 & none \\
Qwen3 & \texttt{k8}, unhealed & 0.107796 & 0.109185 & $-0.139$ & 0.162 & 0.0964 & 0.2185 & none \\
Qwen3 & intact & 0.461104 & 0.112212 & $+34.889$ & 0.844 & 0.0001 & 0.0005 & item signal \\
OLMo-2 & intact & 0.271858 & 0.112324 & $+15.953$ & 0.761 & 0.0001 & 0.0005 & item signal \\
OLMo-2 & \texttt{keep10}@83500 & 0.112450 & 0.111508 & $+0.094$ & 0.481 & 0.6966 & 0.8623 & none \\
OLMo-2 & \texttt{keep12}@124000 & 0.113115 & 0.112219 & $+0.090$ & 0.471 & 0.7026 & 0.8623 & none \\
OLMo-2 & \texttt{keep14}@200000 & 0.119847 & 0.113766 & $+0.608$ & 0.500 & 0.0172 & 0.0663 & trace \\
OLMo-2 & \texttt{shortgpt16}@200000 & 0.153341 & 0.112797 & $+4.054$ & 0.575 & 0.0001 & 0.0005 & item signal \\
Qwen3 & \texttt{k10}, unhealed & 0.118185 & 0.116155 & $+0.203$ & 0.228 & 0.0804 & 0.2185 & none \\
Qwen3 & \texttt{k12}, unhealed & 0.115027 & 0.114880 & $+0.015$ & 0.379 & 0.9432 & 0.9968 & none \\
Qwen3 & \texttt{k14}, unhealed & 0.118933 & 0.116263 & $+0.267$ & 0.188 & 0.0066 & 0.0297 & trace \\
Llama-2 & intact & 0.131981 & 0.112392 & $+1.959$ & 0.525 & 0.0001 & 0.0005 & anchor block \\
Llama-2 & \texttt{k8}, unhealed & 0.107463 & 0.111620 & $-0.416$ & 0.495 & 0.1002 & 0.2185 & none \\
Llama-2 & \texttt{k10}, unhealed & 0.113198 & 0.114575 & $-0.138$ & 0.356 & 0.4510 & 0.7163 & none \\
Llama-2 & \texttt{k12}, unhealed & 0.115858 & 0.116528 & $-0.067$ & 0.081 & 0.1052 & 0.2185 & none \\
Llama-2 & \texttt{k14}, unhealed & 0.116772 & 0.116408 & $+0.036$ & 0.085 & 0.3950 & 0.6666 & none \\
Llama-3 & intact & 0.329205 & 0.112014 & $+21.719$ & 0.813 & 0.0001 & 0.0005 & item signal \\
Llama-3 & \texttt{k8}, unhealed & 0.113531 & 0.115399 & $-0.187$ & 0.351 & 0.2894 & 0.5581 & none \\
Llama-3 & \texttt{k10}, unhealed & 0.111951 & 0.112269 & $-0.032$ & 0.366 & 0.8746 & 0.9839 & none \\
Llama-3 & \texttt{k12}, unhealed & 0.111868 & 0.111851 & $+0.002$ & 0.514 & 0.9968 & 0.9968 & none \\
Llama-3 & \texttt{k14}, unhealed & 0.114694 & 0.110945 & $+0.375$ & 0.446 & 0.0970 & 0.2185 & none \\
\bottomrule
\end{tabular}
\caption{Complete 27-cell v2 read-out. Null is $\widehat{\mathrm{acc}}$;
deltas and CI half-widths are percentage points. ``Item signal,'' ``trace,''
and ``anchor block'' abbreviate the verdicts defined in the text. The final
column reports Benjamini--Hochberg $q$-values over exactly these 27 cells: six
clear $q\leq0.05$ and five clear Bonferroni at $\alpha/27=0.001852$.
Evidence \textsf{E-D}; emitted by \texttt{gate\_bh\_27cell\_consistency.py}.}
\label{tab:v2-full}
\end{table}

% Added columns for ruling P0-5 (materiality constant + Delta_max + recovery_fraction),
% aligned with existing tab:v2-full 27-cell content; values emitted from frozen JSON.
\begin{table}[H]
\centering
\small
\begin{tabular}{@{}lrrrrr@{}}
\toprule
Relative-recovery cutoff $c$ & 0.05 & 0.10 & 0.15 & 0.20 & 0.25 \\
Material cells (all) & 7 & 4 & 4 & 4 & 4 \\
Material cells (damaged only) & 4 & 1 & 1 & 1 & 1 \\
\bottomrule
\end{tabular}

\vspace{4pt}
\footnotesize
\setlength{\tabcolsep}{4.2pt}
\begin{tabular}{lrrrr}
\toprule
Cell & $\Delta_{\mathrm{perm}}$ & $\Delta_{\max}$ & \texttt{recovery\_fraction} & family anchor rf \\
\midrule
qwen3/k8+f2 heal@5000 & -0.072 & 8.372 & -0.0086 & 0.5441 \\
qwen3/k8+f2 heal@5500 & +0.004 & 11.249 & 0.0003 & 0.5441 \\
qwen3/k8+f2 heal@6000 & -0.100 & 12.126 & -0.0082 & 0.5441 \\
qwen3/k8+f2 heal@6500 & -0.092 & 11.062 & -0.0083 & 0.5441 \\
qwen3/k8+f2 heal@7000 & -0.022 & 9.137 & -0.0024 & 0.5441 \\
olmo2/keep8+f2 heal@45000 & +0.322 & 20.718 & 0.0156 & 0.3045 \\
olmo2/keep8+f2 heal@121000 [P2] & +0.229 & 24.955 & 0.0092 & 0.3045 \\
qwen3/k8 UN-healed [P1] & -0.139 & 2.172 & -0.0640 & 0.5441 \\
qwen3 INTACT & +34.889 & 64.120 & 0.5441 & 0.5441 \\
olmo2 INTACT & +15.953 & 52.398 & 0.3045 & 0.3045 \\
olmo2/keep10+f2 heal@83500 & +0.094 & 30.206 & 0.0031 & 0.3045 \\
olmo2/keep12+f2 heal@124000 & +0.090 & 28.763 & 0.0031 & 0.3045 \\
olmo2/keep14+f2 heal@200000 & +0.608 & 29.797 & 0.0204 & 0.3045 \\
olmo2/shortgpt16 heal@200000 & +4.054 & 40.715 & 0.0996 & 0.3045 \\
qwen3/k10 UN-healed & +0.203 & 8.090 & 0.0251 & 0.5441 \\
qwen3/k12 UN-healed & +0.015 & 19.596 & 0.0007 & 0.5441 \\
qwen3/k14 UN-healed & +0.267 & 5.412 & 0.0493 & 0.5441 \\
llama2 INTACT & +1.959 & 35.927 & 0.0545 & 0.0545 \\
llama2/k8 UN-healed & -0.416 & 31.050 & -0.0134 & 0.0545 \\
llama2/k10 UN-healed & -0.138 & 18.172 & -0.0076 & 0.0545 \\
llama2/k12 UN-healed & -0.067 & 1.180 & -0.0568 & 0.0545 \\
llama2/k14 UN-healed & +0.036 & 0.892 & 0.0408 & 0.0545 \\
llama3 INTACT & +21.719 & 66.176 & 0.3282 & 0.3282 \\
llama3/k8 UN-healed & -0.187 & 12.089 & -0.0155 & 0.3282 \\
llama3/k10 UN-healed & -0.032 & 19.982 & -0.0016 & 0.3282 \\
llama3/k12 UN-healed & +0.002 & 33.180 & 0.0001 & 0.3282 \\
llama3/k14 UN-healed & +0.375 & 26.929 & 0.0139 & 0.3282 \\
\bottomrule
\end{tabular}
\caption{\textbf{The material set is stable over the intended cutoff range.} The top panel reports the threshold sweep over the relative-recovery cutoff $c$ (material $\Leftrightarrow \Delta_{\mathrm{perm}}/\Delta_{\max} \ge c \times r_f(\mathrm{intact, same\ family})$); the damaged-only row excludes the three intact anchors (whose own ratios are $1.0$ by construction), so the pre-registered $c=0.10$ operating point retains exactly one damaged cell (\texttt{olmo2/shortgpt16}). the lower panel gives per-cell diagnostics for the 27 cells of Table~\ref{tab:v2-full}: signed $\Delta_{\mathrm{perm}}$ (percentage points), the re-assignment capacity bound $\Delta_{\max}$, the relative-recovery fraction $\Delta_{\mathrm{perm}}/\Delta_{\max}$, and the same-family intact-anchor fraction used by the anchor-competence precondition. \texttt{qwen3/intact}, \texttt{olmo2/intact}, and \texttt{llama3/intact} are their families' anchors. Numbers are emitted from the frozen within-stratum permutation record by \texttt{code/emit\_materiality\_sweep.py}.}
\label{tab:materiality-diagnostics}
\end{table}


\FloatBarrier

\subsection{Rejected v2 alternatives}

\begin{table}[!htbp]
\centering
\scriptsize
\setlength{\tabcolsep}{3.0pt}
\begin{tabular}{p{0.15\linewidth}p{0.17\linewidth}p{0.17\linewidth}p{0.15\linewidth}p{0.26\linewidth}}
\toprule
Alternative & Cell & Null/statistic (pp) & Comparator (pp) & Measured reason for rejection \\
\midrule
Own-modal null & Qwen3 \texttt{k8} & $+1.230$ vs always-E & v2 $-0.139$ & Over-credits by 1.37 pp and flips the sign \\
Content interface & Qwen3 heal@7000 & content $-3.610$ & letter $-0.022$ & More negative under the same permutation null \\
Content interface & Qwen3 \texttt{k8} & content $-1.076$ & letter $-0.139$ & More negative under the same permutation null \\
Content interface & OLMo-2 \texttt{keep8}@121000 & content $-0.086$ & letter $+0.229$ & Does not rescue the damaged read-out \\
Entropy/modal share & Llama-3 \texttt{k12} & $\Delta_{\mathrm{perm}}=+0.002$ & $p=0.997$ & Modal share 0.339 and entropy 0.614 still contain no signal \\
\bottomrule
\end{tabular}
\caption{Measured reasons for rejecting simpler v2 alternatives. Numbers are the same-cell comparisons recorded alongside the adopted permutation null (evidence \textsf{E-D}; see the reproducibility statement for the artifact path).}
\label{tab:v2-alternatives}
\end{table}


\FloatBarrier

\clearpage
\subsection{Integrity repairs}

\begin{table}[!htbp]
\centering
\small
\setlength{\tabcolsep}{4.2pt}
\begin{tabular}{>{\raggedright\arraybackslash}p{0.16\linewidth}>{\raggedright\arraybackslash}p{0.22\linewidth}>{\raggedright\arraybackslash}p{0.28\linewidth}>{\raggedright\arraybackslash}p{0.24\linewidth}}
\toprule
Defect & Consequence & Repair/hardening & Measured impact \\
\midrule
Bootstrap zero atom counted in both tails & Illegal $p=1.042$ & Two-sided mid-$p$ splits zero mass evenly; $p\leq1$ structurally & 0/24 verdicts changed; 0/30 crossed 0.05; non-$p$ fields byte-identical \\
Tokenizer-specific left truncation & 10/15 first-launch cells lost part of labelled option body & Raise cap to 2048; \texttt{n\_trunc} becomes a hard per-shard assertion & 0/14 verdicts changed; maximum 0.0083 pp; all eight argmax flips on affected items, 0 elsewhere \\
Llama-2 KV-cache OOM & 5/8 shards failed; partial 3/8 result unavailable & Merge guard refused output; \texttt{use\_cache=False} for teacher-forced pass & Peak memory 94 GiB to 41--50 GiB; intact cell recovered \\
\bottomrule
\end{tabular}
\caption{Integrity defects and repairs. Numbers are from the bootstrap mid-$p$ repair record and the truncation/OOM re-launch audit of the first cross-family MMLU-Pro run (evidence \textsf{E-E} and \textsf{E-I}; see the reproducibility statement for the artifact paths).}
\label{tab:integrity}
\end{table}


\FloatBarrier

\subsection{Generated construct-null manifest}
\label{app:generated-nulls}

Table~\ref{tab:app-construct-nulls} restates the construct nulls of
Table~\ref{tab:nulls} as a \emph{generated} artefact, and additionally retains the
superseded legality-blind $p$ column beside the corrected one. It is written by
the manifest emitter directly from the balanced-null calibration record
(evidence \textsf{E-CAL}), the single
machine-readable source for these nine rows, and the emitter contains no numeric
literal for any floor or chance line. The emitter refuses to write on any
inconsistency: it checks that each stored percentage-point gap equals
$100\,(\text{Floor}-\text{Chance})$, that each floor is a label count divided by
$n$, that no floor falls below $1/k$, that each row's balanced-null verdict
agrees both with its $p$-value and with the file's own partition of the
constructs, that the corrected record agrees with the record it supersedes on
every quantity the correction does not touch --- $n$, $k$, floor, chance and gap
--- so that a ``correction'' cannot silently move a floor, and that the row
identities agree with a second, independently recorded evidence file. It also
re-parses every number in the hand-authored
Table~\ref{tab:nulls} and in its own output and re-verifies each against the
source at the precision that table chose. A companion checker extends the same
tolerance-based comparison to the prose. Both exit non-zero on any mismatch, so a
typed number that drifts from the evidence becomes a build failure rather than a
reading error. Both are named with their repository paths in
Table~\ref{tab:artifact-map}.

% GENERATED FILE -- DO NOT EDIT BY HAND.
% Emitted by paperC/code/emit_tab_construct_nulls.py from
% paperC/evidence/construct_nulls_legality_aware.json (sha256 ec58acbeb433),
% which supersedes evidence/floor_winners_curse_calibration.json (sha256 275112623d05).
% Regenerate with:  python paperC/code/emit_tab_construct_nulls.py
\begin{table}[H]
\centering
\scriptsize
\setlength{\tabcolsep}{2.6pt}
\begin{tabular}{lrrrrrrrl}
\toprule
 & & & & & & \multicolumn{2}{c}{$p$ under null} & \\
\cmidrule(lr){7-8}
Construct & $n$ & $k$ & Chance & Floor & Gap (pp) & blind & legal & Verdict \\
\midrule
MMLU-Pro letter, naive & 12032 & 10 & 0.100000 & 0.116606 & +1.661 & $<10^{-5}$ & 0.083 & inside estimator noise \\
MMLU-Pro letter, item-avg. & 12032 & 10 & 0.110877 & 0.116606 & +0.573 & $<10^{-5}$ & 0.083 & inside estimator noise \\
MMLU letter & 14042 & 4 & 0.250000 & 0.268908 & +1.891 & $<10^{-5}$ & $<10^{-5}$ & above balanced null \\
OpenBookQA letter & 500 & 4 & 0.250000 & 0.276000 & +2.600 & 0.383 & 0.384 & inside estimator noise \\
ARC-Easy letter & 2376 & 4 & 0.250161 & 0.266414 & +1.625 & 0.140 & 0.136 & inside estimator noise \\
ARC-Challenge letter & 1172 & 4 & 0.250156 & 0.265358 & +1.520 & 0.453 & 0.448 & inside estimator noise \\
CommonsenseQA letter & 1221 & 5 & 0.200000 & 0.208845 & +0.884 & 0.853 & 0.853 & inside estimator noise \\
PIQA letter & 1838 & 2 & 0.500000 & 0.504897 & +0.490 & 0.658 & 0.691 & inside estimator noise \\
BoolQ & 3270 & 2 & 0.500000 & 0.621713 & +12.171 & $<10^{-5}$ & $<10^{-5}$ & above balanced null \\
\bottomrule
\end{tabular}
\caption{\textbf{Construct-null manifest, generated from evidence \textsf{E-CAL}.} The legality-blind column is retained beside the admissible legal null: for the 4 variable-option rows, only the legal null respects each item's available labels. Verdicts therefore follow the legal column. All values, including $k$, Gap, and both $p$ columns, are emitted and checked from the frozen calibration record (schema~1.0.0; seed~20260814; $1000000$ draws; SHA prefix \texttt{ec58acbeb433}); Appendix~\ref{app:generated-nulls} gives the construction and supersession details.}
\label{tab:app-construct-nulls}
\end{table}


\FloatBarrier

\subsection{What we claim and what we retract}
\label{app:ledger}

\begin{table}[H]
\centering
\scriptsize
\setlength{\tabcolsep}{3.5pt}
\begin{tabular}{p{0.05\linewidth}p{0.37\linewidth}p{0.19\linewidth}p{0.31\linewidth}}
\toprule
\# & Claim previously considered & Current status & Manuscript-safe form \\
\midrule
1 & Letter behavior is a family-general step function or sharp phase transition. & Retracted & Report per-family descriptions only. \\
2 & Damage turns the letter interface into a constant predictor. & Narrowed & Damage drives the letter score to or below its best-constant floor; modal share is separate. \\
3 & Exact ties are the family-general causal mechanism. & Retracted & Precision changes ties and argmaxes but not accuracy; mechanisms are family-specific. \\
4 & Letter MC is generally an unreliable instrument. & Retracted & Intact models can favor letter or content; calibrate each reported construct. \\
5 & Accuracy-versus-normalized-accuracy length sensitivity is our finding. & Preempted & Attribute generic over-correction and tokenizer dependence to Oostermeijer; our result is floor inheritance. \\
6 & Content is within $\pm3$ pp of letter on every damaged arm. & Re-scoped & False off MMLU; ARC-Easy reaches $+38.76$ pp. \\
7 & Damaged letter scores are significantly below floor across small non-MMLU benchmarks. & Retracted & Only ARC-Easy \texttt{keep10} in OLMo-2; accompany other nulls with power. \\
8 & \texttt{k14} is the last arm above floor in every family. & Retracted & Descriptive only for healed OLMo-2; other ladders are cliff-like and regime-confounded. \\
9 & A pooled floor verdict can be read per benchmark. & Prohibited & Pool only per-item deviations and label the aggregate scope. \\
10 & A content floor is a dataset property. & Narrowed & Token floors require dataset, convention, unit, and tokenizer; letter floors are dataset properties. \\
11 & Larger BPE vocabularies create more option-length ties. & Retracted & Tokenizer dependence stands; the vocabulary-size mechanism is non-monotone. \\
12 & Significant below-floor behavior is MMLU-specific. & Retracted and reinterpreted & MMLU-Pro has v1 below-floor cells, but v2 shows no negative item-level signal. \\
13 & Damage universally drives letter to or below floor. & Narrowed & The honest MMLU-Pro count is 15/17; Qwen3 \texttt{k14} is a real but immaterial exception and \texttt{shortgpt16} clears its floor decisively. \\
14 & First-launch cross-family MMLU-Pro numbers are usable. & Retracted & Use only the 2048-token corrected launch and canonical v2 evidence. \\
15 & Clearing the floor certifies validity. & Retracted & Clearing one reference is not sufficient; failure withholds a label, success does not certify. \\
16 & Deprecated numeric claims remain usable. & Prohibited & Do not resurrect 4.8$\times$, 0.2822, 58/91, MMLU 0.25, BoolQ 0.50, or 45.74 pp. \\
\bottomrule
\end{tabular}
\caption{Flat current-status ledger. The 16 rows and their status are the retraction/narrowing ledger maintained with the claim-to-evidence map, whose retracted-claim rows correspond one-to-one with the rows here (evidence \textsf{E-K}; see the reproducibility statement for the artifact path).}
\label{tab:claims}
\end{table}


\FloatBarrier

\subsection{Reporting checklist}

For each reported multiple-choice construct, we recommend the following minimal record:
\begin{enumerate}[leftmargin=*]
    \item Define the input-blind predictor family and report the best member on the evaluated item set.
    \item Report the score, null value, paired effect, confidence interval, and test.
    \item For content length nulls, print tie convention, length unit, and tokenizer. If option count varies, print the definition of chance.
    \item Establish that the experiment could resolve the effect under discussion; otherwise describe the result as a non-observation.
    \item Keep arm-independent interface validity separate from arm-conditional item-level information.
    \item Treat a floor pass as necessary, not sufficient.
\end{enumerate}

\subsection{Evidence provenance}
\label{app:provenance}

The construct floors, MMLU-Pro cell table, read-out v2 values, fp32 comparison, small-benchmark power analysis, integrity repairs, and claim ledger are transcribed from a single machine-extracted evidence record, whose lettered sections are the identifiers \textsf{E-A} through \textsf{E-K} used in the table captions. The character-versus-token table additionally uses that record's August 14 addendum, which records the raw-parquet reconstruction and 12/12 self-test. No healed Qwen3 step-121000 result is included.

Table~\ref{tab:artifact-map} resolves every evidence identifier used in a caption to the
file that carries it in the released artifact, together with the emitters and checkers named
above. Captions cite the identifier rather than the path so that the claim remains legible
without the artifact in hand; this table is the one place where the identifiers become
literal paths, so a reader who does have the artifact can follow any caption to its source.
Paths are relative to the accompanying code repository.

\begin{table}[!htbp]
\centering
\small
\setlength{\tabcolsep}{4pt}
\begin{tabular}{lp{0.62\linewidth}}
\toprule
Identifier & Artifact path (relative to the code repository root) \\
\midrule
\textsf{E-A} & \texttt{evidence/SECOND\_MC\_BENCHMARK\_VERDICT.md} \\
\textsf{E-B} & per-cell records in \texttt{evidence/mmlu\_scale\_power/} and \texttt{evidence/POWER\_WALL\_VERDICT.md} \\
\textsf{E-D} & \texttt{evidence/heal\_readout\_v2\_permutation\_null.json} \\
\textsf{E-E} & \texttt{evidence/R7\_BOOTSTRAP\_P\_FIX.md} \\
\textsf{E-F} & \texttt{evidence/second\_mc\_benchmark/} and \texttt{evidence/SECOND\_MC\_BENCHMARK\_VERDICT.md} \\
\textsf{E-H} & \texttt{evidence/construct\_nulls\_length\_unit.json} \\
\textsf{E-I} & \texttt{evidence/mmlu\_scale\_} \texttt{power/trunc\_fix\_} \texttt{before\_after.json} and \texttt{evidence/mmlu\_scale\_} \texttt{power/} \texttt{trunc\_audit\_1536.json}, emitted by \texttt{code/mmlu\_pro\_} \texttt{trunc\_fix\_compare.py} and \texttt{code/mmlu\_pro\_} \texttt{trunc\_audit.py} \\
\textsf{E-K} & row-by-row map in \texttt{evidence/claim\_evidence\_map.tsv} \\
\textsf{E-CAL} & \texttt{evidence/construct\_nulls\_legality\_aware.json}, which supersedes \texttt{evidence/} \texttt{floor\_winners\_} \texttt{curse\_calibration.json} and names it, with its \texttt{sha256}, in its own \texttt{supersedes} field \\
Emitter & \texttt{code/emit\_tab\_construct\_nulls.py} \\
Null recomputation & \texttt{code/recompute\_legality\_aware\_nulls.py} \\
Prose checker & \texttt{code/check\_prose\_vs\_evidence.py} \\
\bottomrule
\end{tabular}
\caption{Evidence identifiers used in the captions, resolved to artifact paths. The identifiers are stable; the paths are where those records live in the released repository.}
\label{tab:artifact-map}
\end{table}
